% Section 5: the pre-registered frozen-set test (E1 + amendment A4).
% Every number is from analysis-output/e1/e1_scores.json, written by e1_analysis.py from
% results/e1/ (evidence commits 370fe66 and ac8a93e). Pre-registration and deviations:
% claims/PREREG_E1_frozen_eval.md; verdicts: claims/E1_RESULTS.md.
\section{A pre-registered test on held-out questions}
\label{sec:confirm}

The sweep in Section~\ref{sec:frontier} was exploratory. Its settings were chosen after
earlier results were seen, every method was fit on the same pool its 250 evaluation
prompts came from, each method drew its own evaluation prompts, and intervals were
computed over responses although AccessEval's pairs are nested: each question appears
with several disability categories, and the dataset lists every question twice. Before
generating anything further we wrote down predictions, falsifiers and stopping gates, and
ran them on one frozen evaluation set (Appendix~\ref{app:e1}).

\paragraph{Design.} The frozen set holds five disability variants of each of 50
questions, 250 items. Every method is fit only on the 633 distinct pairs of the other 96
questions; as a leak control, the projection and CAA are also fit on the full pool. Each
baseline runs at its published setting, the projection at $\alpha{=}1$ with $k{=}40$.
Quality is scored by two judges from different model families, Llama-3.1-8B-Instruct and
Qwen2.5-7B-Instruct, with the same rubric; they agree per item at Spearman
$\rho{=}0.64$. Intervals are bootstrapped over the 50 questions. Three unsteered seeds
agree to within $0.013$ in mean medicalization, and the degeneracy flags are those of
Section~\ref{sec:frontier}.

\begin{table}[t]
\centering\footnotesize
\setlength{\tabcolsep}{3.2pt}
\begin{tabular}{@{}lrrrr@{}}
\toprule
& \multicolumn{1}{c}{$d$ \scriptsize[95\% CI]} & $\Delta$Med & \multicolumn{2}{c}{$\Delta$Quality} \\
\cmidrule(l){4-5}
& & \% & Llama & Qwen \\
\midrule
\multicolumn{5}{@{}l}{\textit{Prompting only}} \\
Minimal prompt & $-0.04$ \scriptsize$[-.14,.06]$ & $-8$ & $-0.03$ & $-0.03$ \\
Explicit prompt & $0.06$ \scriptsize$[-.02,.15]$ & $13$ & $+0.01$ & $+0.07$ \\
\midrule
\multicolumn{5}{@{}l}{\textit{Steering, intact output}} \\
CAD projection & $0.34$ \scriptsize$[.21,.48]$ & $66$ & $\mathbf{-0.59}$ & $-0.23$ \\
CAD mean-diff.\ $\alpha{=}7$ & $0.43$ \scriptsize$[.32,.56]$ & $91$ & $-0.72$ & $-0.24$ \\
CAA & $0.37$ \scriptsize$[.24,.51]$ & $77$ & $-0.80$ & $-0.30$ \\
FairSteer & $0.28$ \scriptsize$[.15,.42]$ & $59$ & $-1.31$ & $-0.61$ \\
LEACE, rank 40 & $0.17$ \scriptsize$[.05,.29]$ & $37$ & $-0.43$ & $-0.23$ \\
\midrule
\multicolumn{5}{@{}l}{\textit{Prompt and steering combined}} \\
Explicit + projection & $\mathbf{0.51}$ \scriptsize$[.37,.66]$ & $\mathbf{97}$ & $-0.70$ & $\mathbf{-0.20}$ \\
\midrule
\multicolumn{5}{@{}l}{\textit{Degenerate output}} \\
SADI $s{=}10$ & $0.30$ \scriptsize$[.10,.53]$ & $49$ & $-3.93$ & $-2.38$ \\
Angular $150^\circ$ & $0.74$ \scriptsize$[.58,.95]$ & $143$ & $-2.07$ & $-0.66$ \\
\bottomrule
\end{tabular}
\caption{The frozen-set test: 250 items from 50 held-out questions, one shared
unsteered reference (medicalization $0.597$; quality $7.85$ Llama, $6.70$ Qwen).
$\Delta$Med above 100 means the model was pushed past neutral. Bold marks the best value
among intact arms that remove at least half the medicalization. Every arm is in Appendix~\ref{app:e1}.}
\label{tab:e1}
\end{table}

\begin{figure*}[t]
\centering
\includegraphics[width=\textwidth]{figure-02-e1-frozen-set.pdf}
\caption{Every arm of the frozen-set test under both judges. Horizontal: debiasing; vertical:
change in judged quality against the unsteered reference. Lines join strengths or ranks of
one operator. Hollow markers are flagged degenerate from the text alone. Error bars are
$95\%$ intervals over the 50 questions. The two judges differ in scale but agree on the
ordering, and the highest-$d$ point in both panels is broken output.}
\label{fig:e1}
\end{figure*}

\paragraph{The highest scores still belong to broken output.} Angular Steering has the
largest effect size of any arm, $d{=}0.74$, by pushing medicalization past neutral to
$-0.26$, and it loses $2.07$ points of judged quality. SADI at its published setting
scores exactly zero on $58\%$ of responses, and a fifth of its responses cannot be parsed
by the judge. Ordered by quality loss, the five published methods fall in the same order
under both judges, as pre-registered: SADI, Angular, FairSteer, CAA, and the projection.
Ordered by $d$, Angular comes first and SADI ranks above FairSteer.

\paragraph{What did not replicate.} The sweep suggested that the projection loses less
quality than CAA at comparable debiasing. On held-out questions the paired difference is
$+0.22$ $[-0.12, 0.57]$ under the Llama judge and $+0.06$ $[-0.03, 0.15]$ under Qwen, and
debiasing is $0.34$ against $0.37$. The point estimates favour the projection and the
intervals include zero, so we treat the two as equally intact at matched debiasing. The
earlier significant result (Appendix~\ref{app:earlier}) computed intervals over
responses, treating variants of one question as independent, which is the error
Section~\ref{sec:geometry} finds in spectral estimates. The leak did not matter: fitting
on the full pool, evaluation questions included, raises $d$ by $0.004$ for the projection
and $0.001$ for CAA.

\subsection{Prompting does not remove the bias, and combines with steering}
\label{sec:prompting}

Prompting outperforms representation steering on AxBench \citep{wu2025axbench}, so we
added two system prompts, worded before any prompted output existed. The minimal one asks
the model to mention a disability only where it changes the answer; the explicit one
tells it not to assume a disabled person needs medical treatment, diagnosis or therapy.
Neither costs quality, and neither removes much bias: the minimal prompt leaves
medicalization where it was ($d{=}{-}0.04$) and the explicit one removes $13\%$
($d{=}0.06$). The projection removes more than the explicit prompt by $0.28$
$[0.14, 0.44]$. Concept removal is where \citet{macocco2026tradeoff} also find prompting
weakest.

The two combine. With the explicit prompt in place, the projection reaches $d{=}0.51$
$[0.37, 0.66]$ and removes $97\%$ of medicalization with no flagged output, $0.08$ to
$0.26$ more than the projection alone, at a quality cost of $0.70$ (Llama) and $0.20$
(Qwen). No intact arm reaches a higher $d$ (Figure~\ref{fig:e1}). The directions were fit
under the default prompt and still act under the new one.

\subsection{What rank does}
\label{sec:rank}

Removing more directions removes more medicalization and costs more quality
(Figure~\ref{fig:e1}, dark blue line). The projection goes from $d{=}0.02$ at rank 1 to
$0.34$ at rank 40 and $0.45$ at rank 80, while the quality change under the Llama judge
goes from $-0.03$ to $-0.59$ to $-0.92$. Rank-$k$ LEACE \citep{belrose2023leace}, erasing
the same top-$k$ coordinates with whitening, shows the same trend more weakly: $0.01$,
$0.17$ $[0.05, 0.29]$ and $0.30$. Its rank-1 and rank-40 intervals overlap by $0.03$, so
the pre-registered prediction that rank 40 beats rank 1 for both operators fails on
LEACE. Removing the single direction of a supervised probe gives $d{=}0.08$
$[-0.01, 0.18]$: one well-chosen direction, removed, does not reproduce the effect of
forty.

What survives is narrower than a claim about the dimensionality of the bias. A single
removed direction does little, and beyond that rank behaves like a strength setting with
the same trade-off every other strength parameter shows. Section~\ref{sec:geometry}
explains why the rank cannot instead be read off the data.
